\documentclass[a4paper,12pt]{ctexart}

\usepackage[margin=2.54cm]{geometry}
\usepackage{newtxmath}
\usepackage{fontspec}
\usepackage{amsmath}
\usepackage{bm}
\usepackage{indentfirst}

% Western text/numerals use Times New Roman and mathematics uses a
% Times-compatible font. Chinese text and headings follow ctexart defaults.
\setmainfont{Times New Roman}

\setlength{\parindent}{2em}
\setlength{\parskip}{0pt}
\linespread{1.25}
\numberwithin{equation}{section}

\newcommand{\vect}[1]{\bm{#1}}
\newcommand{\dd}{\mathop{}\!\mathrm{d}}
\newcommand{\FT}{\mathcal{F}}

\begin{document}

\section{理论}

\subsection{Maxwell方程到波动方程（Wave equation）}
\begin{equation}
  \begin{aligned}
    \nabla\cdot\vect{D}  & = \rho,                                         \\
    \nabla\cdot\vect{B}  & = 0,                                            \\
    \nabla\times\vect{E} & = -\frac{\partial\vect{B}}{\partial t},         \\
    \nabla\times\vect{H} & = \vect{J}+\frac{\partial\vect{D}}{\partial t}.
  \end{aligned}
  \label{eq:Maxwell}
\end{equation}
方程~\eqref{eq:Maxwell} 中所有的场振幅和相位依赖于空间
$\vect{r}=(x,y)$、时间 $t$ 和传播方向 $z$。电位移矢量的表达式为：
\begin{equation}
  \hat{\vect{D}}(\vect{r},\omega,z)
  =\varepsilon_0\varepsilon(\omega)\hat{\vect{E}}(\vect{r},\omega,z)
  +\hat{\vect{P}}(\vect{r},\omega,z).
  \label{eq:Displacement}
\end{equation}
对方程~\eqref{eq:Maxwell} 中的第三式取旋度，并代入
方程~\eqref{eq:Displacement} 化简得：
\begin{equation}
  \nabla^2\vect{E}-\nabla(\nabla\cdot\vect{E})
  -\frac{1}{c^2}\frac{\partial^2}{\partial t^2}
  \int_{-\infty}^{t}\varepsilon(t-t')\vect{E}(\vect{r},t',z)\dd t'
  =\mu_0\left(\frac{\partial\vect{J}}{\partial t}
  +\frac{\partial^2\vect{P}}{\partial t^2}\right).
  \label{eq:time-domain}
\end{equation}
方程~\eqref{eq:time-domain} 中 $\vect{E}$、$\vect{J}$、$\vect{P}$
依赖于 $(\vect{r},t,z)$，变换到空间--频率域：
\begin{equation}
  \nabla^2\hat{\vect{E}}-\nabla(\nabla\cdot\hat{\vect{E}})
  +\frac{\omega^2n^2(\omega)}{c^2}\hat{\vect{E}}
  =\mu_0\left(-\mathrm{i}\omega\hat{\vect{J}}-\omega^2\hat{\vect{P}}\right).
  \label{eq:frequency-domain}
\end{equation}
方程~\eqref{eq:frequency-domain} 中 $\hat{\vect{E}}$、$\hat{\vect{J}}$、
$\hat{\vect{P}}$ 依赖于 $(\vect{r},\omega,z)$，介电常数
$\varepsilon(\omega)=n^2(\omega)$，表示材料的线性折射率。

\subsection{标量波动方程}
假设电场强度垂直于传播方向线性偏振，偏振方向为 $\vect{e}_s$。那么，
$\vect{E}=E\vect{e}_s$，$\vect{P}=P\vect{e}_s$，$\vect{J}=J\vect{e}_s$。
方程~\eqref{eq:frequency-domain} 中的
$\nabla(\nabla\cdot\hat{\vect{E}})$ 因此可以被忽略。方程可写成标量形式：
\begin{equation}
  \left(\partial_z^2+\nabla_\perp^2\right)\hat{E}(\vect{r},\omega,z)
  +k^2(\omega)\hat{E}(\vect{r},\omega,z)
  =-\mu_0\omega^2\hat{P}(\vect{r},\omega,z).
  \label{eq:scalar-form}
\end{equation}
其中，$k(\omega)=n(\omega)\omega/c$。在需要加入电流面密度 $\hat{J}$ 时，
可以将极化强度 $\hat{P}$ 替换为 $\hat{P}+\mathrm{i}\hat{J}/\omega$。

\subsection{一些变换}

\subsubsection{从实验室参考系变换到脉冲参考系}
实验室参考系 $(z,t)$ 变换到脉冲参考系 $(\zeta,\tau)$：
\begin{equation}
  \begin{aligned}
    \zeta & = z,       \\
    \tau  & = t-z/v_g.
  \end{aligned}
\end{equation}
对于某一依赖关系
$\partial_z A(z,t)=\partial_\zeta A(\zeta,\tau)\,\partial_z\zeta
  +\partial_\tau A(\zeta,\tau)\,\partial_z\tau$，因此：
\begin{equation}
  \begin{aligned}
    \partial_z & = \partial_\zeta-\frac{1}{v_g}\partial_\tau, \\
    \partial_t & = \partial_\tau.
  \end{aligned}
  \label{eq:frame}
\end{equation}
通常选脉冲的群速度作为 $v_g$。在频域中，
$\partial_z=\partial_\zeta+\mathrm{i}\omega/v_g$。

\subsubsection{傅里叶变换}

\paragraph{时间傅里叶变换}
时间 $t$ 和频率 $\omega$ 的傅里叶变换为：
\begin{equation}
  \hat{A}(\omega)=\int_{-\infty}^{+\infty}A(t)\exp(\mathrm{i}\omega t)\dd t
  \approx\frac{T}{N}\sum_{n=0}^{N-1}A(t_n)\exp(\mathrm{i}\omega t_n).
\end{equation}
逆变换为：
\begin{equation}
  A(t)=\frac{1}{2\pi}\int_{-\infty}^{+\infty}\hat{A}(\omega)
  \exp(-\mathrm{i}\omega t)\dd\omega
  \approx\frac{1}{T}\sum_{n=0}^{N-1}\hat{A}(\omega_n)
  \exp(-\mathrm{i}\omega_n t).
\end{equation}
其中，$T$ 为时间窗口，$N$ 为采样点数。

对于连续时间傅里叶变换的卷积定理：
\begin{equation}
  \FT[f(t)*g(t)]=\FT[f(t)]\FT[g(t)].
\end{equation}
其中，$\FT$ 表示连续时间傅里叶变换。对于归一化的 DFT（numpy 或者 matlab）：
\begin{equation}
  \FT_D[f(t)*g(t)]=T\FT_D[f(t)]\FT_D[g(t)].
\end{equation}
对于非归一化的 DFT（FFTW 或者 cuFFT）：
\begin{equation}
  \FT_D[f(t)*g(t)]=\frac{T}{N}\FT_D[f(t)]\FT_D[g(t)].
\end{equation}
由于 DFT 默认时间零点从索引 0 开始，一个函数的时间零点要从索引 0 开始，
另一个函数的时间零点在序列中心。

\paragraph{空间二维傅里叶变换}
空间 $(x,y)$ 和空频域 $(k_x,k_y)$ 的傅里叶变换为：
\begin{equation}
  \begin{aligned}
    \hat{A}(k_x,k_y)
     & =\iint_{-\infty}^{+\infty}A(x,y)
    \exp[-\mathrm{i}(k_xx+k_yy)]\dd x\dd y \\
     & \approx\frac{X}{M}\frac{Y}{N}
    \sum_{m=0}^{M-1}\sum_{n=0}^{N-1}A(x_m,y_n)
    \exp[-\mathrm{i}(k_xx_m+k_yy_n)].
  \end{aligned}
\end{equation}
逆变换为：
\begin{equation}
  \begin{aligned}
    A(x,y)
     & =\frac{1}{(2\pi)^2}\iint_{-\infty}^{+\infty}\hat{A}(k_x,k_y)
    \exp[\mathrm{i}(k_xx+k_yy)]\dd k_x\dd k_y                       \\
     & \approx\frac{1}{XY}\sum_{m=0}^{M-1}\sum_{n=0}^{N-1}
    \hat{A}(k_{x_m},k_{y_n})
    \exp[\mathrm{i}(k_{x_m}x+k_{y_n}y)].
  \end{aligned}
\end{equation}
其中，$X$ 和 $Y$ 为两个方向的窗口长度，$M$ 和 $N$ 为两个方向的采样点数。
数值实现里注意 fft 和 ifft 的使用。

\subsection{载波分辨的单向传输方程}

\subsubsection{标量波动方程分解}
\label{sec:factor-method}
采用分解方程~\eqref{eq:scalar-form} 的方式推导 Forward Maxwell Equation（FME）
\cite{Feit:88}，将传播算子分解为前向和后向：
\begin{equation}
  (\partial_z+\mathrm{i}k(\omega))(\partial_z-\mathrm{i}k(\omega))\hat{E}
  =-\nabla_\perp^2\hat{E}-\mu_0\omega^2\hat{P}(\vect{r},\omega,z).
  \label{eq:factor}
\end{equation}
若方程~\eqref{eq:factor} 不含右手边表示衍射和非线性极化的项，方程的解为
前向传输的波与反向传输的波的叠加：
\begin{equation}
  \hat{E}(\omega,z)=\hat{A}_+(\omega)\exp[\mathrm{i}k(\omega)z]
  +\hat{A}_-(\omega)\exp[-\mathrm{i}k(\omega)z].
\end{equation}
假设前向传输的分量远远强于后向传输的分量，
$|\hat{A}_-|\ll|\hat{A}_+|$，那么近似
$\partial_z+\mathrm{i}k(\omega)\approx2\mathrm{i}k(\omega)$ 成立，
方程~\eqref{eq:factor} 近似为 FME：
\begin{equation}
  \frac{\partial\hat{E}}{\partial z}
  =\mathrm{i}k(\omega)\hat{E}
  +\frac{\mathrm{i}}{2k(\omega)}\nabla_\perp^2\hat{E}
  +\frac{\mathrm{i}}{n(\omega)}\frac{\omega}{c}
  \frac{\hat{P}}{2\varepsilon_0}.
  \label{eq:FME}
\end{equation}

利用方程~\eqref{eq:frame} 将方程~\eqref{eq:FME} 变换到脉冲参考系，
利用傅里叶变换，得到 FME 在空频--时频域 $(\vect{k}_\perp,\omega)$ 的表示：
\begin{equation}
  \frac{\partial\tilde{E}}{\partial\zeta}
  =\mathrm{i}\left[k(\omega)-\frac{k_\perp^2}{2k(\omega)}
    -\frac{\omega}{v_g}\right]\tilde{E}
  +\frac{\mathrm{i}\omega}{cn(\omega)}\frac{\tilde{P}}{2\varepsilon_0}.
  \label{eq:FME-local}
\end{equation}

\subsubsection{标准化形式}
在空频--时频域，不仅仅是 FME，任何单向传输方程（Unidirectional Propagation
Equations）都可以化为标准形式：
\begin{equation}
  \frac{\partial\tilde{E}}{\partial z}
  =\mathrm{i}K_z(\vect{k}_\perp,\omega)\tilde{E}
  +\mathrm{i}Q(\vect{k}_\perp,\omega)\frac{\tilde{P}}{2\varepsilon_0}.
  \label{eq:canonical}
\end{equation}
在脉冲坐标下，$Q$ 的形式不变，$K_z$ 变换为 $K_z-\omega/v_g$：
\begin{equation}
  \frac{\partial\tilde{E}}{\partial\zeta}
  =\mathrm{i}\left[K_z(\vect{k}_\perp,\omega)-\frac{\omega}{v_g}\right]\tilde{E}
  +\mathrm{i}Q(\vect{k}_\perp,\omega)\frac{\tilde{P}}{2\varepsilon_0}.
  \label{eq:canonical-local}
\end{equation}
例如，FME 的 $K_z$ 和 $Q$ 分别为：
\begin{equation}
  \begin{aligned}
    K_z^{(\mathrm{FME})}(\vect{k}_\perp,\omega)
     & =k(\omega)-\frac{k_\perp^2}{2k(\omega)}, \\
    Q^{(\mathrm{FME})}(\vect{k}_\perp,\omega)
     & =\frac{\omega}{cn(\omega)}.
  \end{aligned}
\end{equation}

\subsubsection{Slowly Evolving Wave Approximation}
将方程~\eqref{eq:frame} 直接代入方程~\eqref{eq:scalar-form} 并变换到时频域：
\begin{equation}
  \frac{\partial^2\hat{E}}{\partial\zeta^2}
  +\frac{2\mathrm{i}\omega}{v_g}\frac{\partial\hat{E}}{\partial\zeta}
  =-\nabla_\perp^2\hat{E}
  -\left[k^2(\omega)-\frac{\omega^2}{v_g^2}\right]\hat{E}
  -\frac{\omega^2}{c^2}\frac{\hat{P}}{\varepsilon_0}.
  \label{eq:scalar-form-local}
\end{equation}
引入 Slowly Evolving Wave Approximation（SEWA），忽略
方程~\eqref{eq:scalar-form-local} 左手边 $\hat{E}$ 关于 $\zeta$ 的二阶导，即认为
\begin{equation}
  \left|\frac{\partial^2\hat{E}}{\partial\zeta^2}\right|
  \ll\left|\frac{2\omega}{v_g}\frac{\partial\hat{E}}{\partial\zeta}\right|,
\end{equation}
或者等效地：
\begin{equation}
  \left|\frac{\partial\hat{E}}{\partial\zeta}\right|
  \ll\left|\frac{\omega}{v_g}\hat{E}\right|.
  \label{eq:SEWA}
\end{equation}
方程~\eqref{eq:SEWA} 的物理意义是 $\hat{E}(\vect{r},\omega)$ 发生明显变化的
长度尺寸远远大于 $v_g/\omega$。

将 SEWA 应用于方程~\eqref{eq:scalar-form-local} 得到 Forward Wave Equation（FWE），
实验室坐标下 FWE 的标准形式 $K_z^{(\mathrm{FWE})}(\vect{k}_\perp,\omega)$ 和
$Q^{(\mathrm{FWE})}(\vect{k}_\perp,\omega)$ 为：
\begin{equation}
  \begin{aligned}
    K_z^{(\mathrm{FWE})}(\vect{k}_\perp,\omega)
     & =k(\omega)+\frac{v_g}{2\omega}
    \left[\left(k(\omega)-\frac{\omega}{v_g}\right)^2-k_\perp^2\right], \\
    Q^{(\mathrm{FWE})}(\vect{k}_\perp,\omega)
     & =\frac{\omega v_g}{c^2}.
  \end{aligned}
  \label{eq:FWE}
\end{equation}

比较 FME 和 FWE 的线性和非线性算符，可以得到以下关系式：
\begin{equation}
  \begin{aligned}
    K_z^{(\mathrm{FME})}-K_z^{(\mathrm{FWE})}
     & =-\frac{\omega}{2v_g}\left(1-\frac{n(\omega)v_g}{c}\right)^2
    -\frac{k_\perp^2}{2k(\omega)}
    \left(1-\frac{n(\omega)v_g}{c}\right),                              \\
    Q^{(\mathrm{FME})}-Q^{(\mathrm{FWE})}
     & =\frac{\omega}{n(\omega)c}\left(1-\frac{n(\omega)v_g}{c}\right).
  \end{aligned}
  \label{eq:compare}
\end{equation}

\subsubsection{Unidirectional Pulse Propagation Equation}
将方程~\eqref{eq:scalar-form} 变换到空频--时频域：
\begin{equation}
  \left[\partial_z^2+k^2(\omega)-k_\perp^2\right]\tilde{E}
  =-\mu_0\omega^2\tilde{P}.
\end{equation}
引入 $K_z^{(\mathrm{UPPE})}=\sqrt{k^2(\omega)-k_\perp^2}$，上面方程变成：
\begin{equation}
  (\partial_z+\mathrm{i}K_z)(\partial_z-\mathrm{i}K_z)\tilde{E}
  =-\mu_0\omega^2\tilde{P}.
\end{equation}
通过第~\ref{sec:factor-method} 节中类似的分解方法，可以得到
Unidirectional Pulse Propagation Equation（UPPE）：
\begin{equation}
  \frac{\partial\tilde{E}}{\partial z}
  =\mathrm{i}K_z^{(\mathrm{UPPE})}(\vect{k}_\perp,\omega)\tilde{E}
  +\frac{\mathrm{i}\omega^2}
  {c^2K_z^{(\mathrm{UPPE})}(\vect{k}_\perp,\omega)}
  \frac{\tilde{P}}{2\varepsilon_0}.
  \label{eq:UPPE}
\end{equation}
方程~\eqref{eq:UPPE} 取标准形式时，线性和非线性算符为：
\begin{equation}
  \begin{aligned}
    K_z^{(\mathrm{UPPE})}(\vect{k}_\perp,\omega)
     & =\sqrt{k^2(\omega)-k_\perp^2},                     \\
    Q_z^{(\mathrm{UPPE})}(\vect{k}_\perp,\omega)
     & =\frac{\omega^2}{c^2\sqrt{k^2(\omega)-k_\perp^2}}.
  \end{aligned}
  \label{eq:UPPE-canonical}
\end{equation}
当考虑傍轴近似时，$|k_\perp|\ll|k(\omega)|$，UPPE 可以近似为 FME。以 FME 为主。

\subsection{基于包络的传输方程}
电场强度被描述为脉冲包络 $\mathcal{E}$ 和一个载波频率 $\omega_0$ 的叠加：
\begin{equation}
  E(\vect{r}_\perp,t,z)=\mathcal{E}(\vect{r}_\perp,t,z)
  \exp[\mathrm{i}(k_0z-\omega_0t)].
\end{equation}
在脉冲坐标下：
\begin{equation}
  E(\vect{r}_\perp,\tau,\zeta)=\mathcal{E}(\vect{r}_\perp,\tau,\zeta)
  \exp\left\{\mathrm{i}\left(k_0-\frac{\omega_0}{v_g}\right)\zeta
  -\mathrm{i}\omega_0\tau\right\}.
\end{equation}

在实验室坐标下：
\begin{equation}
  \begin{aligned}
    \partial_zE
     & =\exp[\mathrm{i}k_0z-\mathrm{i}\omega_0t](\partial_z+\mathrm{i}k_0)\mathcal{E},      \\
    \partial_tE
     & =\exp[\mathrm{i}k_0z-\mathrm{i}\omega_0t](\partial_t-\mathrm{i}\omega_0)\mathcal{E}.
  \end{aligned}
  \label{eq:lab-frame}
\end{equation}
在脉冲坐标下：
\begin{equation}
  \begin{aligned}
    \partial_\zeta E
     & =\exp\left\{\mathrm{i}\left(k_0-\frac{\omega_0}{v_g}\right)\zeta
    -\mathrm{i}\omega_0\tau\right\}
    \left[\partial_\zeta+\mathrm{i}\left(k_0-\frac{\omega_0}{v_g}\right)\right]\mathcal{E}, \\
    \partial_\tau E
     & =\exp\left\{\mathrm{i}\left(k_0-\frac{\omega_0}{v_g}\right)\zeta
    -\mathrm{i}\omega_0\tau\right\}
    (\partial_\tau-\mathrm{i}\omega_0)\mathcal{E}.
  \end{aligned}
  \label{eq:pulse-frame}
\end{equation}
包络传输方程也有如下的标准形式：
\begin{equation}
  \frac{\partial\widetilde{\mathcal{E}}}{\partial\zeta}
  =\mathrm{i}\mathcal{K}(\vect{k}_\perp,\Omega)\widetilde{\mathcal{E}}
  +\mathrm{i}\mathcal{Q}(\vect{k}_\perp,\Omega)
  \frac{\widetilde{\mathcal{P}}}{2\varepsilon_0}.
  \label{eq:envelope-canonical}
\end{equation}
其中，$\Omega=\omega-\omega_0$。

\subsubsection{Forward Envelope Equation}
将方程~\eqref{eq:frame} 和方程~\eqref{eq:pulse-frame} 代入
方程~\eqref{eq:canonical} 可得：
\begin{equation}
  \begin{aligned}
    \mathcal{K}(\vect{k}_\perp,\Omega)
     & =K_z(\vect{k}_\perp,\omega=\Omega+\omega_0)
    -\kappa(\omega=\Omega+\omega_0),               \\
    \mathcal{Q}(\vect{k}_\perp,\Omega)
     & =Q(\vect{k}_\perp,\omega=\Omega+\omega_0).
  \end{aligned}
\end{equation}
其中，$\kappa(\omega=\Omega+\omega_0)=k_0+(\omega-\omega_0)/v_g$。

\subsection{至此，主要的文本已经完成，还需要总结和非线性极化强度的具体形式。主要的参考文献\cite{couairon_practitioners_2011}}

\bibliographystyle{unsrt}
\bibliography{参考文献}

\end{document}
